\exnumber{00}
\extitle{UX I - Ex 1 - Entretiens utilisateurs}
\exfiles{ex00/PDF, nomme Name\_FirstName\_DeliverableName\_Date\_VersionNumber}
\exauthorize{Those named in the exercise}
\makeheaderfiles

\textbf{Contexte :}

Welcome to the very first quest of this program.

Your mission, should you choose to accept it, is to find a way to improve athletes' experience related to their hydration when running.

This first project is very special: you are asked to end up with a non-digital solution. As a matter of fact, innovation does not always rhyme with digital. So let's be creative and get out of the box!

\textbf{General recommendations:}

\begin{itemize}
\item Get Out of the Box. Be curious. Follow design trends on YouTube, LinkedIn, in podcasts, on Instagram accounts. Explore the internet when you discover a new concept. Build your own design culture.
\item Document everything you do, take notes, open a medium account to identify and show your deliverables, thoughts and feedback. It will be useful for your portoflio!
\item Trust the design process. One step after another.
\item Design needs to fail. 
\end{itemize}
Failure is even a necessary step, but ideally it should happen before a product is launched, during the prototype and test phases
\begin{itemize}
\item Make connections. Use LinkedIn, your friends, friends of friends to frequently meet designers and ask them about how they practice design in their organization. Who know, some day thay may remember you and offer you a contract!
\item Have fun!!
\end{itemize}

\textbf{Instructions :}

\begin{itemize}
\item Read all the scripts very carefully
\item Extract at least 6 insights and write them down on a PDF file
\begin{itemize}
\item Insights must be either shared by several people or consistent with what other people have already experienced
\item For each insight, name who said it.
\end{itemize}
\item Definition of 'insight': A key learning from what you have seen, read or heard. 
\end{itemize}

\textbf{Livrables :}

\begin{itemize}
\item Create a PDF document and write at least 6 insights and the name of the person(s) who said it.
\item Déposez le PDF sur Github, nommé ainsi : « Name\_FirstName\_DeliverableName\_Date\_VersionNumber ». Exemple : « Doe\_John\_MidFiPrototype\_05242024\_V1 ».
\item Ajoutez un titre dans le document.
\item Soignez l'aspect global du document. Il doit être clair, simple et lisible. Vous pouvez regarder canva.com, sans surcharger le dessin de détails.
\end{itemize}

\textbf{Pistes :}

\begin{itemize}
\item Examples of insights from the scripts (non-exhaustive list):
\begin{itemize}
\item Runs several times a week (Joe, Nathan, Guillaume, Anaelle)
\item Doesn't drink too much right before the race because it hurts my stomach (Nathan, Anaelle)
\item I run races with friends or peers (Joe, Nathan, Guillaume)
\end{itemize}
\end{itemize}

\begin{itemize}
\item The following is the interview script used to conduct the seven user interviews.
\item Interview guide:
\end{itemize}

Hello! My name is Agathe and I am conducting research about athletes and people who practice sports regularly. Hence, I'd be very happy to ask you some questions about the way you run, your habits, and the pain points you may have.

This interview will probably last 20 to 30 minutes.

Do you have any questions before we start?

Run

\begin{itemize}
\item What sports do you practice?
\item How often do you run?
\item When was the last time you went out for a run?
\end{itemize}

Races

\begin{itemize}
\item What was the last race you ran?
\item How was it? Tell me about the experience you had.
\item How did you get prepared?
\item What would you have done differently? Why?
\end{itemize}

Hydration

\begin{itemize}
\item What do you remember about the way you hydrated?
\item How did you manage your hydration?
\item Have you ever found yourself in trouble due to hydration? What happened?
\item What would you never do again? Why?
\end{itemize}

That's it, we went through all the questions I had.

Thank you so much for your time!

I have a last favor to ask you: is there anyone you know who runs frequently and would be available for a talk?

